% from ega4-2-fr.tex
\oldpage[IV-2]{70}

de (I, 5.1.8), $(X_{(k')})_{\mathrm{red}}\otimes_{k'}K$ est isomorphe à
$(X_{(K)})_{\mathrm{red}}$, donc est géométriquement réduit~; il en est alors de
même de $(X_{(k')})_{\mathrm{red}}$ en vertu de (4.6.5, (i)).
\end{proof}

\begin{corollary}[4.6.7]
\label{IV.2.4.6.7-fr}
Si $K$ est une extension de type fini de $k$, il existe une extension radicielle
finie $k'$ de $k$ telle que les corps résiduels de l'anneau semi-local
$K\otimes_k k'$ soient séparables sur $k'$.
\end{corollary}

\begin{proof}
Il suffit d'appliquer (4.6.6) à $X=\operatorname{Spec}(A)$, où $A$ est une
$k$-algèbre de type fini dont $K$ est le corps des fractions.
\end{proof}

\begin{corollary}[4.6.8]
\label{IV.2.4.6.8-fr}
Soit $X$ un $k$-préschéma de type fini. Il existe une extension finie $k'$ de
$k$ telle que $(X_{(k')})_{\mathrm{red}}$ soit géométriquement réduit sur $k'$
et que les composantes irréductibles de $X_{(k')}$ soient géométriquement
irréductibles et les composantes connexes de $X_{(k')}$ géométriquement
connexes.
\end{corollary}

\begin{proof}
Cela résulte aussitôt de (4.6.6), de (4.5.10) et (4.5.15).
\end{proof}

\begin{definition}[4.6.9]
\label{IV.2.4.6.9-fr}
Soient $k$ un corps, $X$ un $k$-préschéma, $x$ un point de $X$. On dit que $X$
est \emph{géométriquement réduit} (ou \emph{séparable}) (resp.
\emph{géométriquement ponctuellement intègre}) au point $x$ sur $k$, si, pour
toute extension $k'$ de $k$ et tout point $x'$ de
$X'=X\otimes_k k'$ au-dessus de $x$, $X'$ est réduit (resp. intègre) en $x'$
(i.e., $\mathcal O_{X',x'}$ est réduit (resp. intègre)). On dit que $X$ est
\emph{géométriquement ponctuellement intègre} si $X$ est géométriquement
ponctuellement intègre en tous ses points, autrement dit si, pour toute
extension $k'$ de $k$, tous les anneaux locaux de $X\otimes_k k'$ sont intègres
(auquel cas on dit aussi que $X\otimes_k k'$ est ponctuellement intègre).
\end{definition}

Notons que, pour que $X$ soit géométriquement réduit sur $k$ (4.6.2), il faut
et il suffit qu'il soit géométriquement réduit sur $k$ en tout point $x\in X$.

\begin{proposition}[4.6.10]
\label{IV.2.4.6.10-fr}
Soient $k$ un corps, $X$ un $k$-préschéma, $k'$ une extension de $k$, $x'$ un
point de $X'=X\otimes_k k'$, $x$ son image dans $X$. Pour que $X$ soit
géométriquement réduit (resp. géométriquement ponctuellement intègre) en $x$
sur $k$, il faut et il suffit que $X'$ le soit en $x'$ sur $k'$. En particulier,
pour que $X$ soit géométriquement ponctuellement intègre sur $k$, il faut et il
suffit que $X'$ soit géométriquement ponctuellement intègre sur $k'$.
\end{proposition}

\begin{proof}
Il suffit de prouver la première assertion. La condition étant évidemment
nécessaire, prouvons qu'elle est suffisante. Supposons donc $X'$
géométriquement réduit (resp. géométriquement ponctuellement intègre) au point
$x'$ sur $k'$ et prouvons que $X$ l'est au point $x$ sur $k$, autrement dit que
pour toute extension $k''$ de $k$ et tout point $x''$ de
$X''=X\otimes_k k''$ au-dessus de $x$, $\mathcal O_{X'',x''}$ est réduit (resp.
intègre). Pour cela, notons qu'il existe un point $z$ de $X'\times_X X''$ qui
se projette en $x'$ et $x''$ (I, 3.4.7). Soit $s$ le point de
$\operatorname{Spec}(k'\otimes_k k'')$ image de $z$ (I, 3.4.9)~; posons
$K=k(s)$, $Z=X\otimes_k K$, de sorte que $K$ peut être considéré comme
extension composée de $k'$ et $k''$, et $z$ comme un point de $Z$ dont l'image
dans $X'$ (resp. $X''$) est $x'$ (resp. $x''$). Comme $X'$ est
géométriquement réduit (resp. géométriquement ponctuellement intègre) au point
$x'$ sur $k'$, l'anneau $\mathcal O_{Z,z}$ est réduit (resp. intègre), et comme
$\mathcal O_{Z,z}$ est un $\mathcal O_{X'',x''}$-module fidèlement plat, il
s'ensuit que $\mathcal O_{X'',x''}$ est réduit (resp. intègre)
$(0_{\mathrm I},6.5.1)$.
\end{proof}

\begin{corollary}[4.6.11]
\label{IV.2.4.6.11-fr}
Supposons $k$ parfait (resp. algébriquement clos). Pour que $X$ soit
géométriquement réduit (resp. géométriquement ponctuellement intègre) au point
$x$ sur $k$, il faut et il suffit que $\mathcal O_{X,x}$ soit réduit (resp.
intègre).
\end{corollary}

\begin{proof}
C'est évidemment nécessaire. Inversement, si cette condition est remplie,
alors, pour toute extension $k'$ de $k$,
$\operatorname{Spec}(\mathcal O_{X,x})\otimes_k k'$ est réduit (4.6.4) (resp.
intègre (4.4.4))~;

